\section{Assumptions and limitations}
\label{sec:limits}
\begin{itemize}[leftmargin=1.5em]
\item \textbf{Gaussian additive noise for all latents.} This is what \cref{thm:ng} requires; for an
      isolated style latent the sink property is only needed trivially, but relaxing the Gaussianity
      for that block is not covered by the theorem. The choice between the ANM and the heteroscedastic
      model is discussed in \cref{app:hnm}.
\item \textbf{Regularity.} $p^{(u)}$ has full support on $\R^{n}$ and is third-order differentiable;
      $g$ is a $\mathrm{C}^2$ diffeomorphism onto its image; faithfulness holds.
\item \textbf{Design.} In camera-trap data not every species occurs at every location or at both times of
      day. \Cref{prop:m1req} gives necessary conditions for any design; sufficiency for a full design is
      expected but not proved (\cref{rem:suff}).
\item \textbf{Coupling through the time of day.} Because $T$ enters both blocks, the content and style
      requirements are not independent; \cref{prop:m1req}(c) is the joint bound.
\item \textbf{Pointwise conditions.} \Cref{as:1,as:2} must hold at every value of the latents, not on
      average.
\item \textbf{Style as a random variable.} Style latents vary within a location and time of day. Effects
	  that are fixed within a location and act on the mixing itself (for example a camera-specific
	  transformation $g(\cdot;\theta_e)$) are outside M1.
\item \textbf{Invariance \labelcref{item:P2} is an assumption.} It can be checked on labelled data, for
      example by testing whether features of a single species and time of day carry location information
      beyond what $Z_s$ explains.
\end{itemize}
